\heading{高等数学A（II）-26春-期末}
\noteinfo{原题为图片回忆版；整理时只对题面作最小润色，使题意更明确。解答尽量使用标准判别法、含参量反常积分定理、Fourier 收敛定理与幂级数逐项积分定理，并补足必要的条件验证。}

\begin{question}[20分]
判断下列级数是否收敛；若收敛，进一步判断是绝对收敛还是条件收敛，并给出详细理由。
\begin{enumerate}
  \item $\displaystyle \sum_{n=1}^{\infty}(-1)^n\left(\sqrt[3]{n+1}-\sqrt[3]{n}\right)$；
  \item $\displaystyle \sum_{n=2}^{\infty}\ln\left(1+\frac{(-1)^n}{n}\right)$。
\end{enumerate}
\begin{solution}
\begin{enumerate}
  \item 记
  \[
    b_n=\sqrt[3]{n+1}-\sqrt[3]{n}.
  \]
  因为函数 $x^{1/3}$ 在 $(0,+\infty)$ 上递增且凹，所以
  \[
    b_n>0,\qquad b_{n+1}\le b_n.
  \]
  又 $b_n\to0$，故由 Leibniz 判别法，
  \[
    \sum_{n=1}^{\infty}(-1)^n b_n
  \]
  收敛。

  再看绝对收敛性。对任意 $N\ge1$，
  \[
    \sum_{n=1}^{N}b_n
    =\sum_{n=1}^{N}\left(\sqrt[3]{n+1}-\sqrt[3]{n}\right)
    =\sqrt[3]{N+1}-1\to+\infty .
  \]
  因而 $\sum b_n$ 发散。故原级数条件收敛。

  \item 先看普通收敛性。对 $k\ge1$，有
  \[
  \begin{aligned}
    &\ln\left(1+\frac1{2k}\right)
    +\ln\left(1-\frac1{2k+1}\right)  \\
    &\qquad
    =\ln\frac{2k+1}{2k}+\ln\frac{2k}{2k+1}=0 .
  \end{aligned}
  \]
  设
  \[
    S_N=\sum_{n=2}^{N}\ln\left(1+\frac{(-1)^n}{n}\right).
  \]
  则
  \[
    S_{2m+1}=0,\qquad
    S_{2m}=\ln\left(1+\frac1{2m}\right)\to0 .
  \]
  因此原级数收敛，且其和为 $0$。

  再看绝对收敛性。令 $t_n=(-1)^n/n$，则 $t_n\to0$，且
  \[
    \lim_{n\to\infty}
    \frac{\left|\ln\left(1+\frac{(-1)^n}{n}\right)\right|}{1/n}
    =
    \lim_{n\to\infty}\left|\frac{\ln(1+t_n)}{t_n}\right|
    =1 .
  \]
  由极限比较判别法，
  \[
    \sum_{n=2}^{\infty}
    \left|\ln\left(1+\frac{(-1)^n}{n}\right)\right|
  \]
  与调和级数同敛散，故发散。所以原级数条件收敛。
\end{enumerate}
\end{solution}
\end{question}

\begin{question}[15分]
证明函数项级数
\[
  \sum_{n=1}^{\infty}(-1)^{n-1}\frac{(1-x)x^n}{1-x^{2n}}
\]
在区间 $(0,1)$ 上一致收敛。
\begin{solution}
令
\[
  u_n(x)=\frac{(1-x)x^n}{1-x^{2n}}
  =\frac{x^n}{1+x+\cdots+x^{2n-1}},
  \qquad 0<x<1.
\]
下面验证一致 Leibniz 判别法的条件。

首先，对任意固定的 $x\in(0,1)$，有
\[
  \frac{u_{n+1}(x)}{u_n(x)}
  =
  \frac{x(1-x^{2n})}{1-x^{2n+2}}.
\]
由于
\[
  1-x^{2n+2}-x(1-x^{2n})
  =(1-x)(1+x^{2n+1})>0,
\]
故
\[
  \frac{u_{n+1}(x)}{u_n(x)}<1.
\]
因此 $u_n(x)$ 关于 $n$ 单调递减。

其次，证明 $u_n(x)\to0$ 关于 $x\in(0,1)$ 一致成立。由算术-几何平均不等式，
\[
\begin{aligned}
  1+x+\cdots+x^{2n-1}
  &\ge
  2n\left(x^{0+1+\cdots+(2n-1)}\right)^{1/(2n)}  \\
  &=2n x^{(2n-1)/2}.
\end{aligned}
\]
于是
\[
  0\le u_n(x)
  \le
  \frac{x^n}{2n x^{(2n-1)/2}}
  =
  \frac{\sqrt{x}}{2n}
  \le \frac1{2n},
  \qquad 0<x<1.
\]
故
\[
  \sup_{0<x<1}u_n(x)\le \frac1{2n}\to0.
\]

综上，$u_n(x)\ge0$，对每个 $x\in(0,1)$ 关于 $n$ 单调递减，且 $u_n(x)\to0$ 在 $(0,1)$ 上一致成立。由一致 Leibniz 判别法，
\[
  \sum_{n=1}^{\infty}(-1)^{n-1}u_n(x)
\]
在 $(0,1)$ 上一致收敛。
\end{solution}
\end{question}

\begin{question}[15分]
设
\[
  \varphi(u)=\int_0^{+\infty}\frac{\arctan x}{x^u(2+x^3)}\,\dd x.
\]
\begin{enumerate}
  \item 求函数 $\varphi(u)$ 的定义域 $I$，即使该反常积分收敛的参数 $u$ 的范围；
  \item 证明 $\varphi(u)$ 在定义域 $I$ 上处处连续。
\end{enumerate}
\begin{solution}
\begin{enumerate}
  \item 只需分别考察 $0$ 与 $+\infty$ 附近的收敛性。

  当 $x\to0^+$ 时，
  \[
    \arctan x\sim x,\qquad 2+x^3\sim2,
  \]
  故
  \[
    \frac{\arctan x}{x^u(2+x^3)}
    \sim \frac12 x^{1-u}.
  \]
  因此积分在 $0$ 附近收敛当且仅当
  \[
    1-u>-1,
    \qquad\text{即}\qquad u<2.
  \]

  当 $x\to+\infty$ 时，
  \[
    \arctan x\to\frac\pi2,\qquad 2+x^3\sim x^3,
  \]
  故
  \[
    \frac{\arctan x}{x^u(2+x^3)}
    \sim \frac\pi2 x^{-u-3}.
  \]
  因此积分在 $+\infty$ 附近收敛当且仅当
  \[
    -u-3<-1,
    \qquad\text{即}\qquad u>-2.
  \]

  综上，
  \[
    I=(-2,2).
  \]

  \item 任取闭区间 $[\alpha,\beta]\subset(-2,2)$。对 $u\in[\alpha,\beta]$ 分段寻找与 $u$ 无关的可积控制函数。

  当 $0<x\le1$ 时，由 $\arctan x\le x$，$2+x^3\ge2$，并且 $x^{-u}\le x^{-\beta}$，得
  \[
    0\le
    \frac{\arctan x}{x^u(2+x^3)}
    \le
    \frac12 x^{1-\beta}.
  \]
  因为 $\beta<2$，所以 $x^{1-\beta}$ 在 $(0,1]$ 上可积。

  当 $x\ge1$ 时，由 $\arctan x\le\pi/2$，$2+x^3\ge x^3$，并且 $x^{-u}\le x^{-\alpha}$，得
  \[
    0\le
    \frac{\arctan x}{x^u(2+x^3)}
    \le
    \frac{\pi}{2}x^{-\alpha-3}.
  \]
  因为 $\alpha>-2$，所以 $x^{-\alpha-3}$ 在 $[1,+\infty)$ 上可积。

  被积函数关于参数 $u$ 连续，且在任意闭区间 $[\alpha,\beta]\subset I$ 上有与 $u$ 无关的可积控制函数。由含参量反常积分连续性定理，$\varphi(u)$ 在 $[\alpha,\beta]$ 上连续。由于这样的闭区间可任取，故 $\varphi(u)$ 在 $I=(-2,2)$ 上处处连续。
\end{enumerate}
\end{solution}
\end{question}

\begin{question}[20分]
考虑函数
\[
  f(x)=\frac{x^2}{4}-\frac{\pi x}{2}+\frac{\pi^2}{6},
  \qquad x\in[0,\pi].
\]
求 $f(x)$ 的 $2\pi$ 周期偶延拓的 Fourier 余弦级数，并利用该 Fourier 级数求
\[
  \sum_{n=1}^{+\infty}\frac1{n^2},
  \qquad
  \sum_{n=0}^{+\infty}\frac{(-1)^n}{(2n+1)^3}
\]
的和。
\begin{solution}
设 Fourier 余弦级数为
\[
  f(x)\sim \frac{a_0}{2}+\sum_{n=1}^{\infty}a_n\cos nx,
\]
其中
\[
  a_0=\frac2\pi\int_0^\pi f(x)\,\dd x,
  \qquad
  a_n=\frac2\pi\int_0^\pi f(x)\cos nx\,\dd x
  \quad(n\ge1).
\]

先算常数项：
\[
\begin{aligned}
  a_0
  &=\frac2\pi\int_0^\pi
  \left(\frac{x^2}{4}-\frac{\pi x}{2}+\frac{\pi^2}{6}\right)\dd x  \\
  &=\frac2\pi\left(\frac{\pi^3}{12}-\frac{\pi^3}{4}+\frac{\pi^3}{6}\right)=0.
\end{aligned}
\]

当 $n\ge1$ 时，
\[
  \int_0^\pi x\cos nx\,\dd x=\frac{(-1)^n-1}{n^2},
\]
并且
\[
  \int_0^\pi x^2\cos nx\,\dd x=\frac{2\pi(-1)^n}{n^2}.
\]
于是
\[
\begin{aligned}
  a_n
  &=\frac2\pi\int_0^\pi
  \left(\frac{x^2}{4}-\frac{\pi x}{2}+\frac{\pi^2}{6}\right)\cos nx\,\dd x  \\
  &=\frac2\pi\left[
    \frac14\cdot\frac{2\pi(-1)^n}{n^2}
    -\frac\pi2\cdot\frac{(-1)^n-1}{n^2}
  \right]  \\
  &=\frac1{n^2}.
\end{aligned}
\]
因此
\[
  f(x)\sim \sum_{n=1}^{\infty}\frac{\cos nx}{n^2}.
\]

$f$ 的偶延拓是连续的分段光滑 $2\pi$ 周期函数。由 Fourier 收敛定理，其 Fourier 级数在每一点收敛到函数值本身。故
\[
  f(x)=\sum_{n=1}^{\infty}\frac{\cos nx}{n^2},
  \qquad 0\le x\le\pi.
\]

令 $x=0$，得
\[
  \sum_{n=1}^{\infty}\frac1{n^2}
  =f(0)=\frac{\pi^2}{6}.
\]

下面求第二个数项级数。由于
\[
  \sum_{n=1}^{\infty}\left|\frac{\cos nx}{n^2}\right|
  \le
  \sum_{n=1}^{\infty}\frac1{n^2},
\]
由 Weierstrass 判别法，级数
\[
  \sum_{n=1}^{\infty}\frac{\cos nx}{n^2}
\]
在 $\mathbb R$ 上一致收敛，故可在任意有限区间上逐项积分。于是
\[
\begin{aligned}
  \sum_{n=1}^{\infty}\frac{\sin nx}{n^3}
  &=
  \int_0^x\sum_{n=1}^{\infty}\frac{\cos nt}{n^2}\,\dd t  \\
  &=
  \int_0^x f(t)\,\dd t  \\
  &=
  \frac{x^3}{12}-\frac{\pi x^2}{4}+\frac{\pi^2x}{6}.
\end{aligned}
\]
令 $x=\pi/2$，得
\[
  \sum_{n=1}^{\infty}\frac{\sin(n\pi/2)}{n^3}
  =\frac{\pi^3}{32}.
\]
而
\[
  \sin\frac{n\pi}{2}
  =
  \begin{cases}
    0, & n\ \text{为偶数},\\
    (-1)^k, & n=2k+1,
  \end{cases}
\]
故
\[
  \sum_{k=0}^{\infty}\frac{(-1)^k}{(2k+1)^3}
  =\frac{\pi^3}{32}.
\]

综上，
\[
  \boxed{\displaystyle \sum_{n=1}^{\infty}\frac1{n^2}=\frac{\pi^2}{6}},
  \qquad
  \boxed{\displaystyle \sum_{n=0}^{\infty}\frac{(-1)^n}{(2n+1)^3}=\frac{\pi^3}{32}}.
\]
\end{solution}
\end{question}

\begin{question}[15分]
求幂级数
\[
  \sum_{n=1}^{+\infty}\frac{2}{4n-1}x^{2n}
\]
的和函数，并说明其收敛区间。
\begin{solution}
先求收敛区间。记
\[
  c_n=\frac{2}{4n-1}.
\]
则原级数是关于 $x$ 的幂级数
\[
  \sum_{n=1}^{\infty}c_n x^{2n}.
\]
由根值判别法，
\[
  \lim_{n\to\infty}\sqrt[2n]{c_n}=1,
\]
故收敛半径为 $1$。当 $|x|<1$ 时收敛；当 $|x|>1$ 时一般项不趋于 $0$，发散；当 $x=\pm1$ 时，
\[
  \sum_{n=1}^{\infty}\frac{2}{4n-1}
\]
与调和级数同敛散，故发散。因此收敛区间为
\[
  (-1,1).
\]

下面求 $|x|<1$ 时的和函数。令
\[
  r=\sqrt{|x|},
  \qquad 0\le r<1.
\]
则 $x^{2n}=r^{4n}$，从而
\[
  S(x)=\sum_{n=1}^{\infty}\frac{2}{4n-1}r^{4n}
  =2r\sum_{n=1}^{\infty}\frac{r^{4n-1}}{4n-1}.
\]
对任意固定的 $0\le r<1$，幂级数
\[
  \sum_{n=1}^{\infty}t^{4n-2}
\]
在 $[0,r]$ 上一致收敛。由幂级数逐项积分定理，
\[
\begin{aligned}
  \sum_{n=1}^{\infty}\frac{r^{4n-1}}{4n-1}
  &=
  \sum_{n=1}^{\infty}\int_0^r t^{4n-2}\,\dd t  \\
  &=
  \int_0^r\sum_{n=1}^{\infty}t^{4n-2}\,\dd t  \\
  &=
  \int_0^r\frac{t^2}{1-t^4}\,\dd t .
\end{aligned}
\]
又
\[
  \frac{t^2}{1-t^4}
  =
  \frac12\left(\frac1{1-t^2}-\frac1{1+t^2}\right),
\]
所以
\[
\begin{aligned}
  \int_0^r\frac{t^2}{1-t^4}\,\dd t
  &=
  \frac12\int_0^r\frac{\dd t}{1-t^2}
  -\frac12\int_0^r\frac{\dd t}{1+t^2}  \\
  &=
  \frac12\operatorname{arctanh}r-\frac12\arctan r.
\end{aligned}
\]
于是
\[
  S(x)
  =
  r\left(\operatorname{arctanh}r-\arctan r\right),
  \qquad r=\sqrt{|x|}.
\]
即
\[
  \boxed{\displaystyle
  S(x)=
  \sqrt{|x|}
  \left(
    \operatorname{arctanh}\sqrt{|x|}
    -\arctan\sqrt{|x|}
  \right),
  \qquad |x|<1 .}
\]
当 $x=0$ 时，右端按极限值理解为 $0$，与原级数一致。
\end{solution}
\end{question}

\begin{question}[15分]
设 $a_n>0$，且级数 $\sum_{n=1}^{+\infty}a_n$ 发散。记
\[
  S_n=a_1+a_2+\cdots+a_n.
\]
证明：
\begin{enumerate}
  \item $\displaystyle \sum_{n=1}^{+\infty}\frac{a_n}{1+n^2a_n}$ 收敛；
  \item $\displaystyle \sum_{n=1}^{+\infty}\frac{a_n}{1+a_n}$ 发散；
  \item $\displaystyle \sum_{n=1}^{+\infty}\frac{a_n}{S_n^2}$ 收敛。
\end{enumerate}
\begin{solution}
\begin{enumerate}
  \item 对每个 $n\ge1$，有
  \[
    0<
    \frac{a_n}{1+n^2a_n}
    \le
    \frac{a_n}{n^2a_n}
    =
    \frac1{n^2}.
  \]
  因为 $\sum n^{-2}$ 收敛，所以由比较判别法，
  \[
    \sum_{n=1}^{+\infty}\frac{a_n}{1+n^2a_n}
  \]
  收敛。

  \item 对任意 $t>0$，有
  \[
    \frac{t}{1+t}\ge \frac12\min\{t,1\}.
  \]
  事实上，若 $0<t\le1$，则
  \[
    \frac{t}{1+t}\ge\frac{t}{2};
  \]
  若 $t\ge1$，则
  \[
    \frac{t}{1+t}\ge\frac12.
  \]

  下面证明
  \[
    \sum_{n=1}^{\infty}\min\{a_n,1\}
  \]
  发散。若 $a_n\ge1$ 对无穷多个 $n$ 成立，则上式含有无穷多个不小于 $1$ 的项，必发散。若 $a_n\ge1$ 只对有限多个 $n$ 成立，则从某一项起 $0<a_n<1$，从而
  \[
    \min\{a_n,1\}=a_n
  \]
  最终成立。此时由于 $\sum a_n$ 发散，仍有
  \[
    \sum_{n=1}^{\infty}\min\{a_n,1\}
  \]
  发散。

  因此由比较判别法，
  \[
    \sum_{n=1}^{+\infty}\frac{a_n}{1+a_n}
  \]
  发散。

  \item 因为 $a_n>0$ 且 $\sum a_n$ 发散，所以
  \[
    S_n\nearrow+\infty.
  \]
  对 $n\ge2$，有 $S_n>S_{n-1}>0$，并且
  \[
\begin{aligned}
    \frac1{S_{n-1}}-\frac1{S_n}
    &=
    \frac{S_n-S_{n-1}}{S_{n-1}S_n}  \\
    &=
    \frac{a_n}{S_{n-1}S_n}
    \ge
    \frac{a_n}{S_n^2}.
\end{aligned}
  \]
  于是对任意 $N\ge2$，
  \[
\begin{aligned}
    \sum_{n=2}^{N}\frac{a_n}{S_n^2}
    &\le
    \sum_{n=2}^{N}
    \left(
      \frac1{S_{n-1}}-\frac1{S_n}
    \right)  \\
    &=
    \frac1{S_1}-\frac1{S_N}
    \le
    \frac1{S_1}.
\end{aligned}
  \]
  因为 $\sum a_n/S_n^2$ 是正项级数，其从第二项开始的部分和有上界，所以
  \[
    \sum_{n=2}^{\infty}\frac{a_n}{S_n^2}
  \]
  收敛。再加上首项有限，故
  \[
    \sum_{n=1}^{\infty}\frac{a_n}{S_n^2}
  \]
  收敛。
\end{enumerate}
\end{solution}
\end{question}
